\section{梯度检验：数值验证}

\subsection{数值梯度}

如何验证你手动推导或实现的梯度是否正确？可以用\textbf{数值梯度}（Numerical Gradient）进行检验：

$$\frac{\partial f}{\partial x} \approx \frac{f(x + h) - f(x - h)}{2h}$$

其中 $h$ 是一个很小的数（通常取 $10^{-4}$）。

\begin{figure}[H]
\centering
\includegraphics[width=0.95\textwidth]{slides-images/slide-082.jpg}
\caption{数值梯度检验：使用双侧差分 $\frac{f(x+h) - f(x-h)}{2h}$ 估计梯度，与解析梯度比较\protect\footnotemark}
\end{figure}
\footnotetext{Slides 第83页。}

\begin{warningbox}{双侧差分 vs 单侧差分}
一定要用\textbf{双侧差分}（centered difference）$\frac{f(x+h) - f(x-h)}{2h}$，而不是单侧差分 $\frac{f(x+h) - f(x)}{h}$。双侧差分的误差是 $O(h^2)$，而单侧差分的误差是 $O(h)$，精度差了一个数量级。
\end{warningbox}

\subsection{为什么不直接用数值梯度训练？}

数值梯度计算简单，为什么不直接用它代替反向传播？因为它\textbf{极其缓慢}：需要对每一个参数单独做一次前向计算。如果模型有 $N$ 个参数，就需要 $2N$ 次前向传播（双侧差分）。而反向传播只需要\textbf{一次前向 + 一次反向}就能得到所有参数的梯度。

\begin{importantbox}{梯度检验的正确用法}
数值梯度仅用于\textbf{验证}你的 backward 实现是否正确，绝不用于训练。检验方法：
\begin{enumerate}
    \item 用反向传播计算梯度 $g_{\text{analytic}}$
    \item 用数值差分估计梯度 $g_{\text{numeric}}$
    \item 检查 $|g_{\text{analytic}} - g_{\text{numeric}}| < \epsilon$（通常 $\epsilon \sim 10^{-5}$）
\end{enumerate}
在 PyTorch 已预实现大部分操作的今天，梯度检验的需求减少了，但在实现\textbf{自定义操作}时仍然是必不可少的验证手段。
\end{importantbox}

\subsection{本章小结}

数值梯度提供了一种``暴力但可靠''的梯度估计方法，是验证解析梯度实现正确性的黄金标准。使用双侧差分可以获得更高的精度。但由于计算量与参数数量成正比，数值梯度只适合验证，不适合训练。

%% ========== 总结与延伸 ==========

